\chapter{Topologia geral}\label{ch:b3:topology}

O volume do segundo ano fez análise em espaços métricos: distâncias,
bolas, sequências. Mas as noções fundamentais --- \omterm{def:b3:topology:continuity}{continuidade},
\omterm{def:b3:topology:compact}{compacidade}, \omterm{def:b3:topology:connected}{conexidade} --- jamais mencionam o valor numérico de uma
distância, apenas a família de \emph{abertos} que ela gera. Este
capítulo toma essa família como objeto primitivo. O ganho não é
generalidade pela generalidade: as construções por quociente (o círculo
como $\R/\Z$, os espaços projetivos), os produtos e as \omterm{def:b3:topology:topology}{topologias} do
tipo fraco da análise funcional simplesmente não são objetos que nascem
métricos. Reconstruímos a \omterm{def:b3:topology:continuity}{continuidade}, tratamos em seguida a
\omterm{def:b3:topology:compact}{compacidade} por coberturas abertas (demonstrando que ela equivale, em
espaços métricos, à definição sequencial do segundo ano) e a \omterm{def:b3:topology:connected}{conexidade}
--- encerrando com o teorema de que nenhuma bijeção \omterm{def:b3:topology:continuity}{contínua} pode
identificar $\R$ com $\R^2$: a \omterm{def:b3:topology:topology}{topologia} sabe distinguir dimensões.

\section{Topologias, abertos, continuidade}

\begin{definition}\label{def:b3:topology:topology}
Uma \emph{topologia}\index{topologia} sobre um conjunto $X$ é uma
família $\mathcal T$ de subconjuntos de $X$ --- chamados \emph{conjuntos
abertos}\index{conjunto aberto} --- tal que: $\varnothing, X \in
\mathcal T$; toda reunião de
abertos é aberta; toda interseção \emph{finita} de abertos é aberta. O
par $(X, \mathcal T)$ é um \emph{espaço topológico}. Os complementares dos abertos
são os \emph{fechados}. Uma \emph{vizinhança}\index{vizinhança} de
$x$ é um conjunto que contém um aberto que contém $x$.
\end{definition}

\begin{example}\label{ex:b3:topology:examples}
(a) Um espaço métrico, com ``aberto'' no sentido do segundo ano
(reuniões de bolas abertas): a \emph{\omterm{def:b3:topology:topology}{topologia} métrica}; métricas
diferentes podem dar a mesma \omterm{def:b3:topology:topology}{topologia} (métricas equivalentes). Um
espaço cuja \omterm{def:b3:topology:topology}{topologia} provém de alguma métrica é \emph{metrizável}.
(b) A \omterm{def:b3:topology:topology}{topologia} \emph{discreta} (todos os subconjuntos) e a \omterm{def:b3:topology:topology}{topologia}
\emph{grosseira} $\{\varnothing, X\}$.
(c) A \omterm{def:b3:topology:topology}{topologia} \emph{cofinita} sobre um conjunto infinito: aberto
$=$ vazio ou cofinito. Não é metrizável, como veremos
(\cref{exo:b3:topology:3}).
(d) Em $\R$, a \omterm{def:b3:topology:topology}{topologia} usual; em $\bar\R = \R \cup
\{\pm\infty\}$, a \omterm{def:b3:topology:topology}{topologia} da ordem gerada
pelas semirretas --- o que faz de ``$x_n \to +\infty$'' um caso de convergência
comum.
\end{example}

\begin{definition}\label{def:b3:topology:interior}
Para $A \subseteq X$: o \emph{interior} $\mathring A$ é o maior aberto contido em $A$ (a
reunião de todos eles); o \emph{fecho}\index{fecho} $\bar A$ é o
menor fechado que contém $A$; a \emph{fronteira} é $\partial A = \bar A
\setminus \mathring A$. $A$ é
\emph{denso}\index{conjunto denso} se $\bar A = X$. Tem-se $x \in \bar A$ se, e somente
se, toda \omterm{def:b3:topology:topology}{vizinhança} de $x$ encontra $A$ (se alguma \omterm{def:b3:topology:topology}{vizinhança} não
encontra $A$, o complementar de seu núcleo aberto é um fechado menor
em torno de $A$; e reciprocamente).
\end{definition}

\begin{definition}\label{def:b3:topology:basis}
Uma \emph{base}\index{base de uma topologia} de $\mathcal T$ é uma família $\mathcal B \subseteq \mathcal T$
tal que todo aberto é uma reunião de membros de $\mathcal B$ (por exemplo, as
bolas abertas em um espaço métrico; os intervalos abertos em $\R$).
Uma família $\mathcal B$ de subconjuntos de $X$ é base de \emph{alguma}
\omterm{def:b3:topology:topology}{topologia} se, e somente se, ela cobre $X$ e, para $B_1, B_2 \in \mathcal B$ e $x \in B_1 \cap B_2$,
existe $B_3 \in \mathcal B$ com $x \in B_3 \subseteq B_1\cap B_2$ --- e então ``as reuniões de membros'' formam uma
\omterm{def:b3:topology:topology}{topologia}, a \omterm{def:b3:topology:topology}{topologia} \emph{gerada} por $\mathcal B$.
\end{definition}

\begin{definition}\label{def:b3:topology:continuity}
$f \colon X \to Y$ é \emph{contínua}\index{aplicação contínua} se $f^{-1}(V)$ é aberto para
todo aberto $V \subseteq Y$ --- equivalentemente, as pré-imagens de fechados são
fechadas; equivalentemente, para todo $x$ e toda \omterm{def:b3:topology:topology}{vizinhança} $V$
de $f(x)$, $f^{-1}(V)$ é \omterm{def:b3:topology:topology}{vizinhança} de $x$ (continuidade \emph{em} cada
$x$). Basta verificar sobre uma \omterm{def:b3:topology:basis}{base} de $Y$. Composições de
aplicações contínuas são contínuas. Um
\emph{homeomorfismo}\index{homeomorfismo} é uma bijeção contínua de
inversa contínua; a \omterm{def:b3:topology:topology}{topologia} estuda as propriedades preservadas por
homeomorfismos.
\end{definition}

\begin{proposition}[Continuidade e fecho; colagem]
\label{prop:b3:topology:gluing}
(a) $f$ é \omterm{def:b3:topology:continuity}{contínua} se, e somente se, $f(\bar A) \subseteq \overline{f(A)}$ para todo $A \subseteq X$.
(b) Se $X = F_1 \cup F_2$ com $F_1, F_2$ fechados e $f\colon X
\to Y$ se restringe \omterm{def:b3:topology:continuity}{continuamente} a cada
$F_i$, então $f$ é \omterm{def:b3:topology:continuity}{contínua}.
\end{proposition}

\begin{proof}
(a) Se $f$ é \omterm{def:b3:topology:continuity}{contínua}: $f^{-1}(\overline{f(A)})$ é fechado e contém $A$, logo contém
$\bar A$. Reciprocamente, aplique o critério a $A = f^{-1}(F)$, com $F$
fechado: $f(\bar A) \subseteq
\overline{f(f^{-1}(F))} \subseteq \bar F = F$, logo $\bar A
\subseteq f^{-1}(F) = A$: as pré-imagens de fechados são fechadas.
(b) Para $F \subseteq Y$ fechado: $f^{-1}(F) =
(f\restriction_{F_1})^{-1}(F) \cup (f\restriction_{F_2})^{-1}
(F)$, uma reunião de dois conjuntos fechados em
$F_1$ e em $F_2$, respectivamente, logo fechados em $X$ (os
$F_i$ são fechados: fechado dentro de fechado é fechado).
\end{proof}

\begin{definition}\label{def:b3:topology:hausdorff}
$X$ é \emph{de Hausdorff}\index{espaço de Hausdorff} (ou
\emph{separado}) se dois pontos distintos quaisquer têm \omterm{def:b3:topology:topology}{vizinhanças}
disjuntas. Os espaços métricos são de Hausdorff (bolas de raio $d(x,y)/2$).
Uma sequência $(x_n)$ \emph{converge} para $x$ se toda \omterm{def:b3:topology:topology}{vizinhança}
de $x$ contém todos os $x_n$ salvo um número finito; em um espaço
de Hausdorff, os limites são únicos (dois limites teriam \omterm{def:b3:topology:topology}{vizinhanças}
disjuntas contendo, cada uma, uma cauda). Em um espaço de Hausdorff, os
pontos --- e portanto os conjuntos finitos --- são fechados.
\end{definition}

\begin{remark}\label{rem:b3:topology:sequences}
Em espaços métricos, as sequências detectam tudo: $x \in \bar A$ se, e somente
se, alguma sequência de $A$ converge para $x$ (tome $x_n \in A \cap
B(x, 1/n)$), e
$f$ é \omterm{def:b3:topology:continuity}{contínua} se, e somente se, é sequencialmente \omterm{def:b3:topology:continuity}{contínua}. Em
espaços gerais, ambas as equivalências falham; os substitutos corretos
das sequências (filtros, redes) pertencem a um curso mais avançado.
Enunciaremos os resultados baseados em sequências em espaços métricos e
os baseados em coberturas em geral --- e demonstraremos sua
equivalência onde ela vale.
\end{remark}

\section{Subespaços, produtos, quocientes}

\begin{definition}\label{def:b3:topology:constructions}
Três maneiras de fabricar espaços novos a partir de antigos:\index{topologia
produto}\index{topologia quociente}
\begin{enumerate}
  \item \emph{Subespaço}: em $A \subseteq X$, os abertos são os $U \cap A$, com $U$
        aberto em $X$ --- a \omterm{def:b3:topology:topology}{topologia} menos fina que torna a
        inclusão \omterm{def:b3:topology:continuity}{contínua}.
  \item \emph{Produto}: em $X \times Y$ (e nos produtos finitos), a \omterm{def:b3:topology:topology}{topologia}
        de \omterm{def:b3:topology:basis}{base} os \emph{blocos abertos} $U \times
        V$; em um produto infinito
        $\prod_i X_i$, a \omterm{def:b3:topology:basis}{base} é formada pelos blocos $\prod U_i$ com $U_i = X_i$ para
        \emph{todos salvo um número finito} de índices $i$ --- a
        \omterm{def:b3:topology:topology}{topologia} menos fina que torna \omterm{def:b3:topology:continuity}{contínua} toda projeção.
  \item \emph{Quociente}: se $\sim$ é uma equivalência em $X$ e
        $\pi\colon X \to X/{\sim}$ é a projeção, declare $V
        \subseteq X/{\sim}$ aberto se, e somente se,
        $\pi^{-1}(V)$ é aberto --- a \omterm{def:b3:topology:topology}{topologia} mais fina que torna $\pi$
        \omterm{def:b3:topology:continuity}{contínua}.
\end{enumerate}
\end{definition}

\begin{proposition}[Propriedades universais]
\label{prop:b3:topology:universal}
(a) $f \colon Z \to X \times Y$ é \omterm{def:b3:topology:continuity}{contínua} se, e somente se, ambas as componentes $\pi_X \circ f$,
$\pi_Y \circ f$ o são (o mesmo vale para produtos arbitrários).
(b) $g \colon X/{\sim} \to Z$ é \omterm{def:b3:topology:continuity}{contínua} se, e somente se, $g \circ \pi
\colon X \to Z$ o é.
\end{proposition}

\begin{proof}
(a) Necessidade: composições. Suficiência: basta verificar as
pré-imagens dos blocos da \omterm{def:b3:topology:basis}{base}: $f^{-1}(U \times V) = (\pi_Xf)^{-1}(U)
\cap (\pi_Yf)^{-1}(V)$, aberto (apenas um número finito
de fatores $\neq X_i$ no caso infinito). (b) Necessidade: composição.
Suficiência: para $W \subseteq Z$ aberto, $\pi^{-1}(g^{-1}(W)) = (g\pi)^{-1}(W)$ é aberto, o que, pela definição
da \omterm{def:b3:topology:constructions}{topologia quociente}, significa que $g^{-1}(W)$ é aberto.
\end{proof}

\begin{example}\label{ex:b3:topology:circle}
O quociente $\R/\Z$ (identificando $x$ e $x + n$) é
\omterm{def:b3:topology:continuity}{homeomorfo} ao círculo $S^1 = \{z \in \C : \abs z = 1\}$: a aplicação $x
\mapsto \eu^{2\iu\pi x}$ passa a uma bijeção
\omterm{def:b3:topology:continuity}{contínua} $\R/\Z \to S^1$ (\cref{prop:b3:topology:universal}(b)); sua inversa é \omterm{def:b3:topology:continuity}{contínua} pelo
argumento de \omterm{def:b3:topology:compact}{compacidade} do~\cref{cor:b3:topology:compacthomeo} abaixo (o~\cref{exo:b3:topology:5}
detalha tudo, inclusive por que $\R/\Z$ é \omterm{def:b3:topology:hausdorff}{de Hausdorff} e \omterm{def:b3:topology:compact}{compacto}).
Do mesmo modo, $[0,1]$ com as extremidades coladas é $S^1$, o
quadrado com os lados opostos colados é o toro, e colar é, enfim, um
teorema, não um desenho.
\end{example}

\section{Compacidade}

\begin{definition}\label{def:b3:topology:compact}
Uma \emph{cobertura aberta} de $X$ é uma família $(U_i)_{i\in I}$ de abertos
com $\bigcup U_i = X$. $X$ é \emph{compacto}\index{espaço compacto} se é \omterm{def:b3:topology:hausdorff}{de
Hausdorff} e toda cobertura aberta admite uma \emph{subcobertura
finita}. Equivalentemente (tomando complementares): toda família de
fechados com a \emph{propriedade da interseção finita} (toda
subfamília finita tem interseção não vazia) tem interseção total não
vazia.
\end{definition}

\begin{theorem}[Primeiras propriedades]\label{thm:b3:topology:compactprops}
Seja $X$ \omterm{def:b3:topology:compact}{compacto}.
\begin{enumerate}
  \item Um subconjunto fechado de $X$ é \omterm{def:b3:topology:compact}{compacto}; um subconjunto
        \omterm{def:b3:topology:compact}{compacto} de um \omterm{def:b3:topology:hausdorff}{espaço de Hausdorff} é fechado.
  \item A imagem \omterm{def:b3:topology:continuity}{contínua} de um \omterm{def:b3:topology:compact}{espaço compacto} em um \omterm{def:b3:topology:hausdorff}{espaço de
        Hausdorff} é compacta. Em particular, uma \omterm{def:b3:topology:continuity}{aplicação contínua}
        $f\colon X
        \to \R$ é limitada e atinge seus extremos.
  \item Uma sequência decrescente de fechados não vazios de $X$ tem
        interseção não vazia.
\end{enumerate}
\end{theorem}

\begin{proof}
(1) Sejam $F \subseteq X$ fechado e $(U_i)$ uma cobertura aberta de $F$
(por abertos de $X$): acrescentar $X \setminus F$ dá uma cobertura aberta de
$X$; uma subcobertura finita, retirado $X \setminus F$, cobre $F$. Que o
subespaço é \omterm{def:b3:topology:hausdorff}{de Hausdorff} é claro. Reciprocamente, sejam $K \subseteq Y$ \omterm{def:b3:topology:compact}{compacto},
$Y$ \omterm{def:b3:topology:hausdorff}{de Hausdorff} e $y \notin K$: para cada $x \in K$ escolha abertos disjuntos
$U_x \ni x$, $V_x \ni y$; um número finito dos $U_x$ cobre $K$, e a
interseção dos $V_x$ correspondentes é uma \omterm{def:b3:topology:topology}{vizinhança} de $y$
disjunta da cobertura de $K$ e, portanto, de $K$: o complementar
de $K$ é aberto.

(2) Se $(V_j)$ cobre $f(X)$, então $(f^{-1}(V_j))$ cobre $X$; uma
subcobertura finita embaixo provém da subcobertura finita em cima. A
imagem é \omterm{def:b3:topology:hausdorff}{de Hausdorff} por ser subespaço. Para $f$ real: $f(X)$ é
\omterm{def:b3:topology:compact}{compacto} em $\R$, logo fechado e limitado (cubra por $(-n, n)$ para a
limitação; fechado por (1)), e um conjunto fechado e limitado contém
seu supremo.

(3) Se $\bigcap F_n = \varnothing$, os abertos $X \setminus F_n$ cobrem $X$; um número finito basta,
de modo que algum $F_{n_0} =
\varnothing$ (a sequência é decrescente) --- contradição.
\end{proof}

\begin{corollary}\label{cor:b3:topology:compacthomeo}
Uma \emph{bijeção} \omterm{def:b3:topology:continuity}{contínua} de um \omterm{def:b3:topology:compact}{espaço compacto} em um \omterm{def:b3:topology:hausdorff}{espaço de
Hausdorff} é um \omterm{def:b3:topology:continuity}{homeomorfismo}.
\end{corollary}

\begin{proof}
A inversa é \omterm{def:b3:topology:continuity}{contínua} se, e somente se, as imagens diretas de fechados
são fechadas; um fechado $F$ é \omterm{def:b3:topology:compact}{compacto} (\cref{thm:b3:topology:compactprops}(1)), sua
imagem é compacta (2) e, portanto, fechada (1) no alvo \omterm{def:b3:topology:hausdorff}{de Hausdorff}.
\end{proof}

\begin{theorem}[Produtos finitos]\label{thm:b3:topology:tychonoff}
Um produto finito de \omterm{def:b3:topology:compact}{espaços compactos} é \omterm{def:b3:topology:compact}{compacto}.\index{teorema de
Tychonoff (caso finito)}
\end{theorem}

\begin{proof}
Basta tratar $X \times Y$. A propriedade \omterm{def:b3:topology:hausdorff}{de Hausdorff} é herdada (separe em uma
coordenada). Seja $(W_i)$ uma cobertura aberta de $X
\times Y$; podemos
supor que os $W_i$ são blocos da \omterm{def:b3:topology:basis}{base} $U_i \times
V_i$ (refine: cada ponto está
em um bloco contido em algum $W_i$; uma subcobertura finita de
blocos fornece uma dos $W_i$). Fixe $x \in X$: a \emph{fatia} $\{x\}\times Y \cong Y$ é
compacta, de modo que um número finito de blocos $U_1\times V_1, \dots, U_k\times V_k$ a cobre, com
$x \in U_j$ para todo $j$; então $U^x = \bigcap_{j \leq k} U_j$ é uma \omterm{def:b3:topology:topology}{vizinhança} aberta de $x$
com $U^x \times Y$ coberto por um número finito de blocos (o \emph{lema do
tubo}\index{lema do tubo}: para $(x', y) \in U^x \times Y$, $y \in V_j$ para algum $j$, pois os
blocos cobriam $\{x\}\times Y$ no nível $y$, e $x' \in U^x \subseteq U_j$, logo $(x', y) \in U_j
\times V_j$). Ora, um
número finito de $U^x$ cobre o \omterm{def:b3:topology:compact}{compacto} $X$; as coleções finitas
de blocos correspondentes cobrem $X \times Y$.
\end{proof}

\begin{theorem}[Compacidade em espaços métricos]
\label{thm:b3:topology:metriccompact}
Para um espaço métrico $(X, d)$, são
equivalentes:\index{compacidade sequencial}
\begin{enumerate}
  \item $X$ é \omterm{def:b3:topology:compact}{compacto} (Borel--Lebesgue);
  \item toda sequência em $X$ tem uma subsequência convergente
        (\omterm{thm:b3:topology:metriccompact}{compacidade sequencial} --- a definição do segundo ano);
  \item $X$ é completo e \emph{totalmente limitado}: para todo
        $\varepsilon > 0$, um número finito de bolas de raio $\varepsilon$ cobre $X$.
\end{enumerate}
\end{theorem}

\begin{proof}
(1)$\Rightarrow$(2): suponha que $(x_n)$ não tenha subsequência
convergente. Então cada $x \in X$ tem uma bola aberta $B_x$ que contém
$x_n$ apenas para um número finito de índices $n$ (do
contrário, uma subsequência convergiria para $x$: tome os raios
$1/k$). Um número finito de $B_x$ cobre $X$, de modo que só
existe um número finito de índices $n$: absurdo.

(2)$\Rightarrow$(3): completude: uma sequência de Cauchy com uma subsequência
convergente converge (segundo ano). Total limitação: se algum $\varepsilon$ não
admite cobertura finita, escolha indutivamente $x_{n+1}$ fora de $\bigcup_{k\leq n} B(x_k, \varepsilon)$: a
sequência satisfaz $d(x_m, x_n) \geq \varepsilon$ para $m \neq n$, e não tem subsequência de Cauchy,
nem convergente.

(3)$\Rightarrow$(1): primeiro, (3) implica (2): dada $(x_n)$, cubra $X$
por um número finito de bolas de raio $1$: uma delas, $B_1$,
contém uma subsequência; cubra $X$ por bolas de raio $1/2$: uma
delas contém uma nova subsequência; itere e diagonalize: a subsequência
diagonal é de Cauchy (dois termos posteriores à etapa $k$ estão em
uma bola comum de raio $2^{-k}$, a menos do habitual $2\cdot$) e,
portanto, converge. Seja agora $(U_i)$ uma cobertura aberta e
suponha que não haja subcobertura finita. \emph{Argumento do número de
Lebesgue}: para cada $n$, alguma bola $B(y_n, 2^{-n})$ não é coberta por um
número finito de $U_i$ --- com efeito, cubra $X$ por um número
finito de bolas de raio $2^{-n}$; se cada uma fosse finitamente
coberta, $X$ também seria. Por (2), uma subsequência $y_{n_k} \to y$; escolha
$i$ com $y \in U_i$ e $r > 0$ com $B(y, r) \subseteq U_i$. Para $k$ grande, $B(y_{n_k}, 2^{-n_k}) \subseteq B(y, r) \subseteq U_i$:
coberta por \emph{um único} $U_i$ --- contradição.
\end{proof}

\begin{corollary}[Heine--Borel; Heine]
\label{cor:b3:topology:heineborel}
(a) Um subconjunto de $\R^n$ é \omterm{def:b3:topology:compact}{compacto} se, e somente se, é fechado
e limitado.
(b) Uma \omterm{def:b3:topology:continuity}{aplicação contínua} de um espaço métrico \omterm{def:b3:topology:compact}{compacto} em um espaço
métrico é uniformemente \omterm{def:b3:topology:continuity}{contínua}.\index{teorema de
Heine--Borel}\index{teorema de Heine}
\end{corollary}

\begin{proof}
(a) Fechado e limitado $\Rightarrow$ contido em um cubo $[-M,M]^n$, que é \omterm{def:b3:topology:compact}{compacto}:
$[-M, M]$ o é (sequencialmente, por Bolzano--Weierstrass --- ou diretamente
por dicotomia para as coberturas), e o~\cref{thm:b3:topology:tychonoff} cuida do produto;
depois aplique o~\cref{thm:b3:topology:compactprops}(1). Reciprocamente, um subconjunto
\omterm{def:b3:topology:compact}{compacto} é fechado (\cref{thm:b3:topology:compactprops}(1)) e limitado (cubra por bolas
concêntricas).

(b) Sejam $f \colon X \to Y$, $\varepsilon > 0$. As bolas $B\bigl(x, \delta_x\bigr)$ com $f\bigl(B(x,
2\delta_x)\bigr)\subseteq B\bigl(f(x), \varepsilon/2\bigr)$ cobrem $X$;
extraia uma subcobertura finita $B(x_j, \delta_{x_j})$ e ponha $\delta = \min_j \delta_{x_j}$. Se $d(x, x') < \delta$: $x
\in B(x_j, \delta_{x_j})$
para algum $j$, e então ambos $x, x'
\in B(x_j, 2\delta_{x_j})$, de modo que $d(f(x), f(x')) < \varepsilon$.
\end{proof}

\begin{definition}\label{def:b3:topology:locallycompact}
$X$ é \emph{localmente compacto}\index{espaço localmente compacto}
se é \omterm{def:b3:topology:hausdorff}{de Hausdorff} e todo ponto tem uma \omterm{def:b3:topology:topology}{vizinhança} compacta ($\R^n$;
os abertos de $\R^n$; os espaços discretos --- mas não $\Q$,
veja o~\cref{exo:b3:topology:9}). Todo espaço localmente \omterm{def:b3:topology:compact}{compacto} mergulha em um
\omterm{def:b3:topology:compact}{compacto}: a \emph{compactificação por um ponto}\index{compactificação
por um ponto} $\hat X = X
\cup \{\infty\}$, cujos abertos são os de $X$ junto com os
complementares (em $\hat X$) dos subconjuntos \omterm{def:b3:topology:compact}{compactos} de $X$.
Os axiomas se verificam diretamente; $\hat X$ é \omterm{def:b3:topology:compact}{compacto} (uma
cobertura tem um membro contendo $\infty$, cujo complementar é
\omterm{def:b3:topology:compact}{compacto}, coberto por um número finito dos demais) e \omterm{def:b3:topology:hausdorff}{de Hausdorff}
(separe $x$ de $\infty$ por uma \omterm{def:b3:topology:topology}{vizinhança} compacta de $x$ e
seu complementar). Exemplo: $\hat\R \cong S^1$, e $\widehat{\R^n} \cong S^n$ por projeção estereográfica
(\cref{exo:b3:topology:11}).
\end{definition}

\section{Conexidade}

\begin{definition}\label{def:b3:topology:connected}
$X$ é \emph{conexo}\index{espaço conexo} se não é a reunião de dois
abertos disjuntos não vazios --- equivalentemente, seus únicos
subconjuntos simultaneamente abertos e fechados são $\varnothing$ e $X$;
equivalentemente, toda \omterm{def:b3:topology:continuity}{aplicação contínua} $X \to \{0, 1\}$ (discreto) é constante.
Um subconjunto é conexo se o for como subespaço.
\end{definition}

\begin{theorem}\label{thm:b3:topology:connectedbasics}
\begin{enumerate}
  \item Os subconjuntos \omterm{def:b3:topology:connected}{conexos} de $\R$ são exatamente os
        intervalos.
  \item As imagens \omterm{def:b3:topology:continuity}{contínuas} de \omterm{def:b3:topology:connected}{conexos} são \omterm{def:b3:topology:connected}{conexas} (donde o teorema
        do valor intermediário: uma aplicação real \omterm{def:b3:topology:continuity}{contínua} em um
        \omterm{def:b3:topology:connected}{espaço conexo} tem por imagem um intervalo).
  \item Se os $(A_i)$ são \omterm{def:b3:topology:connected}{conexos} com um ponto em comum, $\bigcup
        A_i$ é
        \omterm{def:b3:topology:connected}{conexo}. Se $A$ é \omterm{def:b3:topology:connected}{conexo} e $A \subseteq
        B \subseteq \bar A$, então $B$ é \omterm{def:b3:topology:connected}{conexo}.
  \item Produtos finitos de \omterm{def:b3:topology:connected}{conexos} são \omterm{def:b3:topology:connected}{conexos}.
\end{enumerate}
\end{theorem}

\begin{proof}
Ao longo de toda a demonstração usamos o critério de $\{0,1\}$: toda $f \colon X \to
\{0,1\}$
\omterm{def:b3:topology:continuity}{contínua} tem de ser constante.

(1) Um $A$ que não é intervalo deixa de conter algum $c$ entre
$a, b \in A$: $A = (A \cap (-\infty, c)) \sqcup (A \cap (c, +\infty))$ o desconecta. Reciprocamente, sejam $I$ um intervalo e
$f \colon I
\to \{0,1\}$ \omterm{def:b3:topology:continuity}{contínua} com $f(a) = 0$, $f(b) = 1$, $a < b$. Seja $c = \sup\{x \in [a,b] : f(x) = 0\}$; a
\omterm{def:b3:topology:continuity}{continuidade} em $c$ força $f(c) = 0$ (limite de valores $0$: toda
\omterm{def:b3:topology:topology}{vizinhança} de $c$ encontra $\{f = 0\}$; $\{f=0\}$ é fechado) e, então,
$c < b$ com $f \equiv 1$ em $(c, b]$, de modo que $f(c) = 1$ pelo mesmo
argumento de \omterm{def:b3:topology:interior}{fecho} sobre $\{f = 1\} \ni c$: contradição.

(2) Uma \omterm{def:b3:topology:continuity}{aplicação contínua} $g \colon f(X) \to \{0,1\}$ dá $g \circ f$ constante, de modo que
$g$ é constante na imagem.

(3) Uma \omterm{def:b3:topology:continuity}{aplicação contínua} $f\colon \bigcup A_i \to \{0,1\}$ é constante em cada $A_i$, com o
mesmo valor no ponto comum. Quanto ao \omterm{def:b3:topology:interior}{fecho}: $f \colon B \to \{0,1\}$ é constante $= c$
em $A$; todo $x \in B \subseteq \bar A$ está no \omterm{def:b3:topology:interior}{fecho} de $A$, e $f^{-1}
(f(x))$ é uma
\omterm{def:b3:topology:topology}{vizinhança} de $x$ (pré-imagem de um aberto), que tem de encontrar
$A$: $f(x) = c$.

(4) Para $X \times Y$ e $f\colon X\times Y \to \{0,1\}$: dois pontos quaisquer $(x,y), (x',y')$ são ligados pelo
``cotovelo'' $\{x\}\times Y \cup X \times \{y'\}$, reunião de dois \omterm{def:b3:topology:connected}{conexos} (\omterm{def:b3:topology:continuity}{homeomorfos} a $Y$ e a
$X$) que se encontram em $(x, y')$: por (3) e (2), $f$ coincide nos
dois pontos.
\end{proof}

\begin{definition}\label{def:b3:topology:pathconnected}
$X$ é \emph{conexo por caminhos}\index{espaço conexo por caminhos}
se dois pontos quaisquer são ligados por um \emph{caminho} (uma
\omterm{def:b3:topology:continuity}{aplicação contínua} $\gamma
\colon [0,1] \to X$). \omterm{def:b3:topology:connected}{Conexo} por caminhos implica \omterm{def:b3:topology:connected}{conexo}: dois
valores de uma $f \colon X \to \{0,1\}$ \omterm{def:b3:topology:continuity}{contínua} em $x, y$ são valores da aplicação
constante $f\circ\gamma$ (\cref{thm:b3:topology:connectedbasics}(1)--(2)). Os subconjuntos convexos de
espaços normados são \omterm{def:b3:topology:connected}{conexos} por caminhos (segmentos); também o são
$\R^n
\setminus \{0\}$ para $n \geq 2$ (contorne a origem) e $S^n$ para $n \geq 1$ (projete
caminhos a partir de $\R^{n+1}\setminus
\{0\}$).
\end{definition}

\begin{example}[A curva seno do topólogo]
\label{ex:b3:topology:sinecurve}
Sejam $\Gamma = \{(x, \sin\frac1x) : 0 < x \leq 1\}$ e $S =
\bar\Gamma = \Gamma \cup (\{0\}\times[-1,1])$ (todo ponto $(0,
y)$, com $\abs y \leq 1$, é limite de
pontos de $\Gamma$: resolva $\sin\frac1x = y$ perto de $0$). Então $S$ é
\emph{\omterm{def:b3:topology:connected}{conexo}} --- \omterm{def:b3:topology:interior}{fecho} do \omterm{def:b3:topology:connected}{conexo} $\Gamma$, imagem \omterm{def:b3:topology:continuity}{contínua} de
$(0, 1]$ (\cref{thm:b3:topology:connectedbasics}(3)) --- mas \emph{não é \omterm{def:b3:topology:pathconnected}{conexo por
caminhos}}: um \omterm{def:b3:topology:pathconnected}{caminho} de $(1, \sin 1)$ até $(0,0)$ teria de percorrer
abscissas $\to 0$ enquanto a ordenada oscila entre $\pm1$; o~\cref{exo:b3:topology:9} torna isso rigoroso. \omterm{def:b3:topology:connected}{Conexidade} e \omterm{def:b3:topology:connected}{conexidade} por
\omterm{def:b3:topology:pathconnected}{caminhos} diferem de verdade.\index{curva seno do topólogo}
\end{example}

\begin{omfigure}
\begin{tikzpicture}
\begin{axis}[omaxis, width=10.5cm, height=4.6cm,
    xmin=-0.02, xmax=0.55, ymin=-1.25, ymax=1.25,
    xtick={0, 0.25, 0.5}, ytick={-1, 1}]
  \addplot[omThm, thick, domain=0.006:0.55, samples=1200]
    {sin(deg(1/x))};
  \addplot[omDef, very thick] coordinates {(0,-1) (0,1)};
\end{axis}
\end{tikzpicture}

{\small A \omterm{ex:b3:topology:sinecurve}{curva seno do topólogo}: o gráfico de $\sin\frac1x$ (em vermelho)
acumula-se sobre todo o segmento $\{0\}\times[-1,1]$ (em azul). A reunião é \omterm{def:b3:topology:connected}{conexa}
mas não é \omterm{def:b3:topology:pathconnected}{conexa por caminhos}: nenhum \omterm{def:b3:topology:pathconnected}{caminho} \omterm{def:b3:topology:continuity}{contínuo} pode atravessar
as infinitas oscilações em tempo de parâmetro finito.}
\end{omfigure}

\begin{definition}\label{def:b3:topology:components}
A \emph{componente conexa}\index{componente conexa} de $x
\in X$ é a
reunião de todos os \omterm{def:b3:topology:connected}{conexos} que contêm $x$ --- o maior deles
(\cref{thm:b3:topology:connectedbasics}(3)). As componentes formam uma partição de $X$ e são
fechadas (\omterm{def:b3:topology:interior}{fechos} de \omterm{def:b3:topology:connected}{conexos} são \omterm{def:b3:topology:connected}{conexos}). $X$ é \emph{totalmente
desconexo} se todas as componentes são unitárias ($\Q$; o conjunto
de Cantor do problema de fim de semana).
\end{definition}

\begin{theorem}\label{thm:b3:topology:rnotr2}
$\R$ não é \omterm{def:b3:topology:continuity}{homeomorfo} a nenhum $\R^n$ com $n \geq 2$.
\end{theorem}

\begin{proof}
Suponha que $h \colon \R \to \R^n$ seja um \omterm{def:b3:topology:continuity}{homeomorfismo}. Retirando um ponto: $\R \setminus \{0\}$ é
\omterm{def:b3:topology:continuity}{homeomorfo} a $\R^n \setminus
\{h(0)\}$. Mas $\R\setminus\{0\}$ é desconexo, ao passo que $\R^n\setminus\{\text{ponto}\}$ é \omterm{def:b3:topology:pathconnected}{conexo
por caminhos} para $n \geq
2$ (\cref{def:b3:topology:pathconnected}; translade o ponto para
$0$): a \omterm{def:b3:topology:connected}{conexidade} é um invariante por \omterm{def:b3:topology:continuity}{homeomorfismo}
(\cref{thm:b3:topology:connectedbasics}(2)) --- contradição. (Que $\R^2 \not\cong \R^3$ exige invariantes mais
finos --- \omterm{def:b3:topology:topology}{topologia} \omterm{def:b3:galois:algebraic}{algébrica}; o problema de fim de semana mostra o
perigo: \emph{sobrejeções \omterm{def:b3:topology:continuity}{contínuas}} $\R \to \R^2$ existem de fato.)
\end{proof}

\begin{method}\label{met:b3:topology:connectargument}
Para demonstrar que um conjunto é \omterm{def:b3:topology:connected}{conexo}: exiba-o como imagem \omterm{def:b3:topology:continuity}{contínua},
como reunião de \omterm{def:b3:topology:connected}{conexos} que se sobrepõem, como \omterm{def:b3:topology:interior}{fecho} ou como produto
(\cref{thm:b3:topology:connectedbasics}); para subconjuntos de espaços normados, demonstre a
\omterm{def:b3:topology:connected}{conexidade} por \omterm{def:b3:topology:pathconnected}{caminhos} com \omterm{def:b3:topology:pathconnected}{caminhos} explícitos (segmentos, arcos,
cotovelos). Para demonstrar que dois espaços \emph{não} são
\omterm{def:b3:topology:continuity}{homeomorfos}: encontre um invariante \omterm{def:b3:topology:topology}{topológico} que difira ---
\omterm{def:b3:topology:compact}{compacidade}, \omterm{def:b3:topology:connected}{conexidade}, número de componentes, ou as componentes
depois de retirado um conjunto finito bem escolhido (o truque
do~\cref{thm:b3:topology:rnotr2}: ele também demonstra $[0,1) \not\cong
(0,1)$ e $S^1 \not\cong [0,1]$).
\end{method}

\section{Exercícios}

\begin{exercise}[$\star$]\label{exo:b3:topology:1}
(a) Liste todas as \omterm{def:b3:topology:topology}{topologias} sobre $\{a, b\}$, classifique-as a menos de
\omterm{def:b3:topology:continuity}{homeomorfismo} e determine quais são \omterm{def:b3:topology:connected}{conexas} e quais são \omterm{def:b3:topology:hausdorff}{de Hausdorff}.
(b) Em $\R$ (\omterm{def:b3:topology:topology}{topologia} usual), calcule \omterm{def:b3:topology:interior}{interior}, \omterm{def:b3:topology:interior}{fecho} e fronteira
de $\Q$, de $[0,1) \cup \{2\}$ e de $\{1/n : n \geq
1\}$.
\end{exercise}

\begin{exercise}[$\star$]\label{exo:b3:topology:2}
(a) Mostre que $f\colon X \to Y$ é \omterm{def:b3:topology:continuity}{contínua} se, e somente se, $f^{-1}(B)$ é aberto para
todo $B$ em uma \omterm{def:b3:topology:basis}{base} fixada de $Y$.
(b) Mostre que $f = \mathbf 1_{[0,\infty)} \colon \R \to \R$ não é \omterm{def:b3:topology:continuity}{contínua}, mas é \emph{\omterm{def:b3:topology:continuity}{contínua} à direita}.
Verifique que os intervalos semiabertos $[a, b)$ formam a \omterm{def:b3:topology:basis}{base de
uma topologia} na fonte $\R$ (a \emph{reta de Sorgenfrey}), que essa
\omterm{def:b3:topology:topology}{topologia} é estritamente mais fina do que a usual, e que uma função
$\R \to
\R$ (com alvo usual) é \omterm{def:b3:topology:continuity}{contínua} a partir da reta de Sorgenfrey se, e
somente se, é \omterm{def:b3:topology:continuity}{contínua} à direita em todo ponto.
\end{exercise}

\begin{exercise}[$\star$]\label{exo:b3:topology:3}
Em um conjunto infinito com a \omterm{def:b3:topology:topology}{topologia} cofinita, mostre: dois abertos
não vazios quaisquer se encontram (de modo que o espaço não é \omterm{def:b3:topology:hausdorff}{de
Hausdorff} e, portanto, não é metrizável); toda sequência injetora
converge para \emph{todo} ponto. Onde a demonstração da unicidade dos
limites usa Hausdorff?
\end{exercise}

\begin{exercise}[$\star\star$]\label{exo:b3:topology:4}
(a) Mostre que as projeções $\pi_X, \pi_Y$ de um produto são \omterm{def:b3:topology:continuity}{contínuas} e
\emph{abertas} (as imagens de abertos são abertas), mas não fechadas em
geral ($\{xy = 1\}$ em $\R^2$).
(b) Mostre que uma sequência em um produto enumerável $\prod_n X_n$ de espaços
métricos converge se, e somente se, cada coordenada converge, e que
$d(x, y) = \sum_n 2^{-n}\min(1, d_n(x_n, y_n))$ metriza a \omterm{def:b3:topology:constructions}{topologia produto}.
(c) Em $\R^\N$, mostre que a ``\omterm{def:b3:topology:topology}{topologia} das caixas'' (todos os
produtos de abertos são abertos) é estritamente mais fina: a sequência
$x^{(k)} = (1/k, 1/k, \dots)$ converge para $0$ na \omterm{def:b3:topology:constructions}{topologia produto}, mas não na \omterm{def:b3:topology:topology}{topologia}
das caixas.
\end{exercise}

\begin{exercise}[$\star\star$]\label{exo:b3:topology:5}
O círculo, de três maneiras. Mostre que os seguintes espaços são dois a
dois \omterm{def:b3:topology:continuity}{homeomorfos}, com aplicações explícitas: (i) $S^1 \subseteq \R^2$; (ii)
$\R/\Z$ (\omterm{def:b3:topology:constructions}{topologia quociente}); (iii) $[0,1]/(0 \sim 1)$. \emph{(Para (ii):
mostre que $\R/\Z$ é \omterm{def:b3:topology:hausdorff}{de Hausdorff} --- levante duas classes a
representantes a distância $\leq \frac12$ --- e que $\pi([0,1])$ é o espaço todo, de
modo que o quociente é \omterm{def:b3:topology:compact}{compacto}; use então o~\cref{cor:b3:topology:compacthomeo}.)}
\end{exercise}

\begin{exercise}[$\star\star$]\label{exo:b3:topology:6}
Seja $X$ um espaço métrico.
(a) Mostre que uma reunião finita de \omterm{def:b3:topology:compact}{compactos} é compacta, e que uma
interseção arbitrária também é.
(b) Se $K$ é \omterm{def:b3:topology:compact}{compacto}, $F$ é fechado e $K \cap F = \varnothing$, mostre que $d(K, F) = \inf\{d(x,y) : x \in K, y \in F\} > 0$;
dê um contraexemplo com dois fechados disjuntos.
(c) Mostre que $X$ é \omterm{def:b3:topology:compact}{compacto} se, e somente se, \emph{toda}
\omterm{def:b3:topology:continuity}{aplicação contínua} $f
\colon X \to \R$ é limitada. \emph{(Se alguma sequência não tem
subsequência convergente, construa uma função \omterm{def:b3:topology:continuity}{contínua} ilimitada com
suporte perto de seus termos; ou use funções do tipo $x \mapsto
d(x, \cdot)$ --- um
\omterm{def:b3:topology:pathconnected}{caminho} limpo: se $(x_n)$ não tem ponto de acumulação, o conjunto
$\{x_n\}$ é fechado e discreto, e $f(x_n) = n$ se estende \omterm{def:b3:topology:continuity}{continuamente} sem
recorrer a Tietze: $f(x) = \sum_n n\,\max\bigl(0, 1 - \frac{d(x, x_n)}{r_n}\bigr)$ para $r_n$ suficientemente pequenos.)}
\end{exercise}

\begin{exercise}[$\star\star$]\label{exo:b3:topology:7}
(Número de Lebesgue) Seja $(U_i)$ uma cobertura aberta de um espaço
métrico \omterm{def:b3:topology:compact}{compacto} $X$. Mostre que existe $\delta > 0$ tal que todo
subconjunto de diâmetro $< \delta$ está contido em um único $U_i$.
\emph{(Do contrário, tome $A_n$ de diâmetro $< 1/n$ em nenhum
$U_i$, e um ponto de acumulação de $x_n \in A_n$ escolhidos.)} Deduza de novo
o teorema de Heine (\cref{cor:b3:topology:heineborel}(b)).
\end{exercise}

\begin{exercise}[$\star\star$]\label{exo:b3:topology:8}
(a) Mostre que $GL_n(\R)$ é um subconjunto aberto e \omterm{def:b3:topology:interior}{denso} de $M_n(\R)$, e que é
desconexo: o sinal do determinante o separa em (ao menos) duas peças.
Mostre, por outro lado, que $GL_n(\C)$ é \omterm{def:b3:topology:pathconnected}{conexo por caminhos}. \emph{(Para
$A, B \in GL_n(\C)$: $\lambda \mapsto \det\bigl((1 -
\lambda)A + \lambda B\bigr)$ é um polinômio não identicamente nulo, de modo que tem um
número finito de raízes em $\C$: escolha um \omterm{def:b3:topology:pathconnected}{caminho} de $\lambda$ em
$\C$ de $0$ a $1$ que as evite.)}
(b) Demonstre a densidade: $A + \varepsilon I$ é inversível para $\varepsilon > 0$ pequeno.
\end{exercise}

\begin{exercise}[$\star\star$]\label{exo:b3:topology:9}
(a) Mostre que as \omterm{def:b3:topology:components}{componentes conexas} de $\Q$ são os conjuntos
unitários, e que $\Q$ não é \omterm{def:b3:topology:locallycompact}{localmente compacto} (uma \omterm{def:b3:topology:topology}{vizinhança}
compacta de $0$ em $\Q$ conteria $[-r, r]\cap\Q$, cujo \omterm{def:b3:topology:interior}{fecho} em $\Q$
não é \omterm{def:b3:topology:compact}{compacto}: corte em um irracional).
(b) Complete o~\cref{ex:b3:topology:sinecurve}: nenhum \omterm{def:b3:topology:pathconnected}{caminho} liga $(0,0)$ a
$\Gamma$. \emph{(Se $\gamma = (\gamma_1, \gamma_2)$ é um tal \omterm{def:b3:topology:pathconnected}{caminho}, seja $t^* = \sup\{t : \gamma_1(t) = 0\}$; para $t
> t^*$ o
ponto se move sobre o gráfico; escolha $t_k \downarrow
t^*$ com $\gamma_1(t_k) \to 0$ atingindo
abscissas em que $\sin$ vale alternadamente $\pm1$ --- o teorema
do valor intermediário as fornece ---, contradizendo a \omterm{def:b3:topology:continuity}{continuidade} de
$\gamma_2$ em $t^*$.)}
\end{exercise}

\begin{exercise}[$\star\star\star$]\label{exo:b3:topology:10}
O conjunto de Cantor dos terços médios $C = \bigcap_n C_n$, em que $C_0 =
[0,1]$ e $C_{n+1}$
remove o terço médio aberto de cada intervalo de $C_n$.
(a) Mostre que $C = \bigl\{\sum_{n\geq1} a_n3^{-n} : a_n \in
\{0,2\}\bigr\}$, e que $C$ é \omterm{def:b3:topology:compact}{compacto}, de \omterm{def:b3:topology:interior}{interior} vazio, e não
tem ponto isolado (é \emph{\omterm{prop:b3:galois:perfect}{perfeito}}).
(b) Mostre que $C$ é totalmente desconexo.
(c) Mostre que $x \mapsto (a_n(x))_n$ é um \omterm{def:b3:topology:continuity}{homeomorfismo} $C \to
\{0,2\}^\N$ (\omterm{def:b3:topology:constructions}{topologia produto} sobre o
espaço discreto de dois pontos); deduza que $C$ é não enumerável e
que $C \times C
\cong C$.
\end{exercise}

\begin{exercise}[$\star\star\star$]\label{exo:b3:topology:11}
Projeção estereográfica: a partir do polo norte $N = (0, \dots,
0, 1)$ de $S^n \subseteq \R^{n+1}$, a
aplicação $\sigma(x) =
\frac{(x_1, \dots, x_n)}{1 - x_{n+1}}$ é um \omterm{def:b3:topology:continuity}{homeomorfismo} $S^n
\setminus \{N\} \to \R^n$ (dê a inversa explicitamente).
Deduza que $S^n$ é \omterm{def:b3:topology:continuity}{homeomorfo} à \omterm{def:b3:topology:locallycompact}{compactificação por um ponto} $\widehat{\R^n}$,
e que retirar \emph{qualquer} ponto de $S^n$ deixa um espaço
\omterm{def:b3:topology:continuity}{homeomorfo} a $\R^n$.
\end{exercise}

\begin{exercise}[$\star\star\star$]\label{exo:b3:topology:12}
(A \omterm{ex:b3:topology:sinecurve}{curva seno do topólogo}) Seja
\[
T = \bigl\{(x, \sin\tfrac1x) : 0 < x \leq 1\bigr\}
\cup \bigl(\{0\}\times\intcc{-1}1\bigr) \subseteq \R^2 .
\]
(a) Mostre que $T$ é \omterm{def:b3:topology:compact}{compacto} e que é o \omterm{def:b3:topology:interior}{fecho} da parte do gráfico.
(b) Mostre que $T$ é \omterm{def:b3:topology:connected}{conexo} \emph{(o gráfico é \omterm{def:b3:topology:connected}{conexo} por ser
imagem \omterm{def:b3:topology:continuity}{contínua}; seu \omterm{def:b3:topology:interior}{fecho} permanece \omterm{def:b3:topology:connected}{conexo})}.
(c) Mostre que $T$ \emph{não} é \omterm{def:b3:topology:pathconnected}{conexo por caminhos}: nenhum \omterm{def:b3:topology:pathconnected}{caminho}
\omterm{def:b3:topology:continuity}{contínuo} liga $(0,0)$ a $(1, \sin1)$. \emph{(Se $\gamma =
(\gamma_1, \gamma_2)$ é um tal \omterm{def:b3:topology:pathconnected}{caminho},
seja $t^* = \sup\{t :
\gamma_1(t) = 0\}$; logo após $t^*$, $\gamma_1$ assume todos os valores
positivos pequenos (teorema do valor intermediário), de modo que $\gamma_2$
oscila entre $\pm1$ em todo intervalo $(t^*, t^* + \delta)$ --- contradiga a
\omterm{def:b3:topology:continuity}{continuidade} em $t^*$.)}
(d) Conclua que a \omterm{def:b3:topology:connected}{conexidade} por \omterm{def:b3:topology:pathconnected}{caminhos} é estritamente mais forte do
que a \omterm{def:b3:topology:connected}{conexidade}, e mostre que nenhum exemplo desse tipo pode ser
aberto em $\R^n$: um subconjunto \emph{aberto} \omterm{def:b3:topology:connected}{conexo} de $\R^n$
é \omterm{def:b3:topology:pathconnected}{conexo por caminhos} \emph{(o conjunto dos pontos que se podem ligar a
um \omterm{def:b3:topology:basis}{ponto-base} por um \omterm{def:b3:topology:pathconnected}{caminho} é aberto e fechado no domínio)}.
\end{exercise}

\section{Problema: o conjunto de Cantor e uma curva que preenche o
quadrado}

\begin{problem}[{Problema de fim de semana --- as curvas de Peano
existem, e por que não são homeomorfismos}]\label{pb:b3:topology:1}
Em 1890, Peano espantou a análise com uma \emph{sobrejeção \omterm{def:b3:topology:continuity}{contínua}}
$[0,1] \to [0,1]^2$: uma curva que preenche um quadrado. Construímos uma delas com as
próprias mãos a partir do conjunto de Cantor $C$ do~\cref{exo:b3:topology:10} (cujos resultados podem ser usados livremente) e, em
seguida, demonstramos que nenhuma tal aplicação pode ser injetora:
quadrados não são curvas. Ao longo de todo o problema, os elementos de
$C$ são escritos $x =
\sum_{n \geq 1} \frac{2b_n(x)}{3^n}$ com dígitos $b_n(x) \in
\{0, 1\}$.

\medskip
\noindent\textbf{Parte I --- Ler dígitos de maneira \omterm{def:b3:topology:continuity}{contínua}.}
\begin{enumerate}
  \item Mostre que, se $x, y \in C$ satisfazem $\abs{x - y} <
        3^{-m}$, então $b_n(x) = b_n(y)$ para todo
        $n \leq m$. \emph{(Se o primeiro dígito diferente é $b_k$, então
        $\abs{x - y} \geq \frac{2}{3^k} - \sum_{j>k}\frac2{3^j}
        = \frac1{3^k}$.)}
  \item Deduza que cada função dígito $b_n \colon C \to \{0,
        1\}$ é \omterm{def:b3:topology:continuity}{contínua} e redemonstre o
        \omterm{def:b3:topology:continuity}{homeomorfismo} $C
        \cong \{0,1\}^\N$ do~\cref{exo:b3:topology:10}(c).
\end{enumerate}

\medskip
\noindent\textbf{Parte II --- Uma sobrejeção \omterm{def:b3:topology:continuity}{contínua} $C \to
[0,1]^2$.}
\begin{enumerate}[resume]
  \item Defina $g \colon C \to [0,1]$ por $g(x) = \sum_{n\geq1}
        b_n(x)\,2^{-n}$ (leia os dígitos de Cantor como
        binários). Mostre que $g$ é \omterm{def:b3:topology:continuity}{contínua} (use a questão 1) e
        sobrejetora. Ela é injetora?
  \item Defina $\Phi \colon C \to [0,1]^2$ por
        \[
        \Phi(x) = \Bigl(\ \sum_{n\geq1} b_{2n-1}(x)\,2^{-n},\
        \sum_{n\geq1} b_{2n}(x)\,2^{-n}\ \Bigr)
        \]
        (os dígitos ímpares dão a abscissa; os pares, a ordenada).
        Mostre que $\Phi$ é \omterm{def:b3:topology:continuity}{contínua} e sobrejetora.
\end{enumerate}

\medskip
\noindent\textbf{Parte III --- Preenchendo as lacunas: a curva de
Peano.}
\begin{enumerate}[resume]
  \item O complementar $[0,1] \setminus C$ é uma reunião enumerável de intervalos
        abertos disjuntos $(u, v)$ (os terços médios removidos)
        cujas extremidades estão em $C$. Defina $f \colon
        [0,1] \to [0,1]^2$ por $f = \Phi$ em
        $C$, estendida \emph{afinamente} em cada lacuna:
        \[
        f(t) = \frac{v - t}{v - u}\,\Phi(u) +
        \frac{t - u}{v - u}\,\Phi(v),
        \qquad t \in (u, v).
        \]
        Mostre que $f$ está bem definida e é sobrejetora sobre
        $[0,1]^2$.
  \item Mostre que $f$ é \omterm{def:b3:topology:continuity}{contínua} em todo ponto de $[0,1]
        \setminus C$
        (localmente afim) e em todo ponto de $C$: dado $\varepsilon$, tome
        $m$ com $2^{-m} <
        \varepsilon/2$ e use a questão 1 para controlar $\Phi$ em
        $C$ perto de $x$; verifique então que a interpolação
        afim não pode escapar: em uma lacuna $(u,v) \subseteq
        (x - 3^{-m}, x + 3^{-m})$, os valores $f(t)$
        estão sobre o segmento $[\Phi(u), \Phi(v)]$, com ambas as extremidades próximas
        de $\Phi(x)$. Conclua: $f$ é uma \emph{sobrejeção \omterm{def:b3:topology:continuity}{contínua}}
        $[0,1] \to [0,1]^2$.
  \item Deduza sobrejeções \omterm{def:b3:topology:continuity}{contínuas} $[0,1] \to [0,1]^n$ para todo $n \geq 2$, e $\R \to \R^2$.
\end{enumerate}

\medskip
\noindent\textbf{Parte IV --- Mas nunca injetora.}
\begin{enumerate}[resume]
  \item Mostre que uma \emph{injeção} \omterm{def:b3:topology:continuity}{contínua} $f \colon [0,1]
        \to [0,1]^2$ seria um
        \omterm{def:b3:topology:continuity}{homeomorfismo} sobre sua imagem (\cref{cor:b3:topology:compacthomeo}).
  \item Mostre que $[0,1]$ e $[0,1]^2$ não são \omterm{def:b3:topology:continuity}{homeomorfos}: retire um
        ponto bem escolhido e compare a \omterm{def:b3:topology:connected}{conexidade} (\cref{met:b3:topology:connectargument}).
  \item Conclua: uma sobrejeção \omterm{def:b3:topology:continuity}{contínua} $[0,1] \to [0,1]^2$ nunca pode ser injetora
        --- uma injetora tornaria $[0,1]^2 = f([0,1])$ \omterm{def:b3:topology:continuity}{homeomorfo} a $[0,1]$ pela
        questão 8, contradizendo a questão 9. Onde exatamente entrou no
        argumento a \omterm{def:b3:topology:compact}{compacidade} de $[0,1]$?
  \item (Culminância) Junte a moral: existe uma sobrejeção \omterm{def:b3:topology:continuity}{contínua},
        mas nenhuma bijeção \omterm{def:b3:topology:continuity}{contínua} $[0,1] \to
        [0,1]^2$. O que isso diz sobre a
        ``dimensão'' como noção \omterm{def:b3:topology:topology}{topológica}? Formule com precisão um
        teorema demonstrado neste problema e um enunciado plausível que
        permanece fora do alcance de nossas ferramentas (invariância do
        domínio).
\end{enumerate}

\medskip
\noindent\textbf{Parte V --- A aritmética do conjunto de Cantor.}
\begin{enumerate}[resume]
  \item Demonstre que $C + C = [0, 2]$: dado $t \in [0, 1]$, escreva $t = \sum_n d_n3^{-n}$ com dígitos $d_n \in \{0,
        1, 2\}$
        e reparta cada dígito como $d_n = a_n + b_n$ com $a_n, b_n \in \{0, 1\}$; conclua que $\frac C2 +
        \frac C2 = [0, 1]$
        \emph{(que subconjunto de $[0,1]$ é $\frac C2$, em termos de
        dígitos ternários?)} e reescale. Nunca é preciso ``vai um'' ---
        diga por que esse é o ponto crucial.
  \item Interprete geometricamente: a ``poeira de Cantor'' $C \times
        C \subseteq \R^2$, de
        \omterm{def:b3:topology:interior}{interior} vazio, projeta uma sombra \emph{cheia} sobre a
        diagonal: a projeção $(x,
        y) \mapsto x + y$ a leva sobre $[0, 2]$. Registre
        também a autossemelhança $C = \frac C3 \cup \bigl(\frac C3 +
        \frac23\bigr)$, a equação por trás de toda
        figura de $C$.
  \item Mostre que o entrelaçamento de dígitos define um \omterm{def:b3:topology:continuity}{homeomorfismo}
        $C
        \to C \times C$ e deduza $C \cong C^n$ para todo $n$ --- o conjunto de Cantor
        é seu próprio quadrado, cubo, \dots{} Que espaços familiares
        compartilham essa propriedade?
  \item Mostre que $\frac14 \in C$ \emph{(calcule sua expansão ternária:
        $\frac14 = \sum_{k\geq1} 2\cdot3^{-2k}$)}, embora $\frac14$ não seja extremidade de nenhum intervalo
        removido; deduza --- contando as extremidades --- que as
        extremidades formam um subconjunto enumerável e próprio de
        $C$.
  \item Mostre que a aplicação de leitura binária $g$ da questão 3
        é no máximo $2$ para $1$ e descreva exatamente quais
        pontos de $[0,1]$ têm duas pré-imagens. ($g$ colapsa
        $C$ sobre $[0,1]$ colando enumeráveis pares: a sombra
        combinatória da escada de Cantor, que reencontraremos no~\cref{ch:b3:measure}.)
\end{enumerate}

\medskip
\noindent\textbf{Parte VI --- \omterm{def:b3:topology:compact}{Compactos} \omterm{prop:b3:galois:perfect}{perfeitos}: Cantor em toda
parte.} Um espaço métrico \omterm{def:b3:topology:compact}{compacto} não vazio $P$ é \emph{\omterm{prop:b3:galois:perfect}{perfeito}}
se não tem ponto isolado.
\begin{enumerate}[resume]
  \item Sejam $P$ \omterm{prop:b3:galois:perfect}{perfeito} \omterm{def:b3:topology:compact}{compacto}, $x \in P$ e $r > 0$. Mostre
        que a bola $B(x, r)$ contém dois pontos $y_0
        \neq y_1$ de $P$ e, portanto,
        duas bolas fechadas \emph{disjuntas} $\bar B(y_0, \rho)$, $\bar B(y_1, \rho)$ dentro de
        $B(x, r)$, centradas em pontos de $P$, de raio $\rho
        \leq r/4$ tão pequeno
        quanto se queira. Explique por que a perfeição (ausência de
        pontos isolados) é exatamente o que permite repetir essa cisão
        dentro de \emph{cada} uma das duas novas bolas.
  \item Itere: construa fechados $F_w \neq \varnothing$ indexados por palavras binárias
        finitas $w$, com $F_{w0},
        F_{w1} \subseteq F_w$ disjuntos e $\operatorname{diam}F_w \leq 2^{-\abs w}
        \operatorname{diam}P$. Mostre que, para
        toda palavra infinita $\varepsilon \in \{0,1\}^\N$, a interseção $\bigcap_mF_{\varepsilon_1\cdots\varepsilon_m}$ é um único ponto
        $h(\varepsilon)$ \emph{(propriedade da interseção finita do \omterm{def:b3:topology:compact}{compacto}
        $P$)}.
  \item Mostre que $h\colon\{0,1\}^\N \to P$ é injetora e \omterm{def:b3:topology:continuity}{contínua}, e conclua: \emph{todo
        espaço métrico \omterm{def:b3:topology:compact}{compacto} \omterm{prop:b3:galois:perfect}{perfeito} é não enumerável} --- na
        verdade, de cardinalidade ao menos a de $\{0,1\}^\N$. Recupere:
        $[0,1]$ e $C$ são não enumeráveis.
  \item Mostre que $h$ é um \omterm{def:b3:topology:continuity}{homeomorfismo} sobre sua imagem
        \emph{(injeção \omterm{def:b3:topology:continuity}{contínua} a partir de um \omterm{def:b3:topology:compact}{compacto},
        \cref{cor:b3:topology:compacthomeo})}: todo espaço métrico \omterm{def:b3:topology:compact}{compacto} \omterm{prop:b3:galois:perfect}{perfeito} contém
        uma cópia \omterm{def:b3:topology:continuity}{homeomorfa} do conjunto de Cantor. O conjunto de
        Cantor não é uma esquisitice, mas o germe universal da
        perfeição compacta.
  \item Deduza que todo espaço métrico \omterm{def:b3:topology:compact}{compacto} enumerável tem um ponto
        isolado, e exiba um em que os pontos isolados são \omterm{def:b3:topology:interior}{densos} mas
        não são tudo: $K = \{0\} \cup
        \{\frac1n : n \geq 1\}$.
  \item (Cantor--Bendixson para $\R$) Seja $F \subseteq \R$ fechado. Diga que
        $x$ é um \emph{ponto de condensação} de $F$ se toda
        \omterm{def:b3:topology:topology}{vizinhança} de $x$ encontra $F$ de maneira não enumerável.
        Mostre que os pontos de condensação de um fechado não
        enumerável $F$ formam um fechado \omterm{prop:b3:galois:perfect}{perfeito} não vazio $P$,
        e que $F
        \setminus P$ é enumerável \emph{(cubra os pontos que não são de
        condensação por enumeráveis intervalos racionais que encontram
        $F$ de maneira enumerável)}. Conclua: todo subconjunto
        fechado de $\R$ é enumerável ou tem a cardinalidade do
        \omterm{def:b3:topology:continuity}{contínuo} --- a hipótese do \omterm{def:b3:topology:continuity}{contínuo} vale para os fechados.
\end{enumerate}

\medskip
\noindent\textbf{Parte VII --- Codas: quão regular, e quão longe da
perfeição.}
\begin{enumerate}[resume]
  \item (Regularidade de Hölder da curva) Ponha $\alpha =
        \frac{\ln 2}{2\ln 3} \approx 0.3155$ e note a
        identidade $3^\alpha = \sqrt2$. Usando a questão 1, mostre que
        \[
        \norm{\Phi(x) - \Phi(y)}_\infty \leq 2\,\abs{x -
        y}^\alpha \qquad (x, y \in C),
        \]
        e propague a estimativa através das lacunas afins: mostre que a
        curva de Peano $f$ da questão 6 satisfaz $\norm{f(t) - f(s)}_\infty \leq 6\abs{t - s}^\alpha$ em todo
        $[0,1]$ \emph{(trate um par na mesma lacuna por
        interpolação, e depois um par geral passando pelos pontos
        extremos de $C \cap [t, s]$)}.
  \item (O expoente $\frac12$ é uma barreira) Mostre que nenhuma sobrejeção
        $F \colon [0,1] \to [0,1]^2$ pode ser $\alpha$-hölderiana com $\alpha > \frac12$: corte
        $[0,1]$ em $N$ intervalos, limite os diâmetros de suas
        imagens e conte os pontos da grade $\bigl(\frac iM, \frac jM\bigr)$, $0 \leq i, j \leq M$, que um conjunto
        de diâmetro $< \frac1M$ pode conter; escolha $M$ da ordem de $N^\alpha$
        e faça $N \to \infty$. Localize nossa curva ($\alpha \approx 0.3155$) em relação à barreira
        e registre, sem demonstração, que a curva de Hilbert atinge o
        expoente crítico $\frac12$.
  \item (Conjuntos derivados: medindo a imperfeição) Para $F$
        fechado em $\R$, seja $F'$ o conjunto dos pontos de
        acumulação de $F$ (um fechado), e itere: $F^{(0)} = F$, $F^{(i+1)}
        = (F^{(i)})'$.
        Verifique que
        \[
        K = \{0\} \cup \Bigl\{\frac1m\Bigr\}_{m \geq 1} \cup
        \Bigl\{\frac1m + \frac1n : m \geq 1,\ n > m^2\Bigr\}
        \]
        é um \omterm{def:b3:topology:compact}{compacto} enumerável com $K' = \{0\} \cup
        \{\frac1m\}$, $K'' = \{0\}$, $K''' = \varnothing$ \emph{(verifique
        que o $m$-ésimo aglomerado vive no intervalo $(\frac1m, \frac1{m-1})$, de modo
        que os aglomerados não se entrelaçam)}, e que seus pontos
        isolados são \omterm{def:b3:topology:interior}{densos} em $K$, como prevê a questão 21.
        Descreva a indução que produz, para todo $j$, um \omterm{def:b3:topology:compact}{compacto}
        enumerável $K_j$ com $K_j^{(j)} = \{0\}$ e $K_j^{(j+1)} =
        \varnothing$, e contraste com os
        conjuntos \omterm{prop:b3:galois:perfect}{perfeitos} da questão 22, para os quais a derivação
        nunca se move: ter posto finito é o exato oposto da perfeição.
\end{enumerate}
\end{problem}
